\documentclass[12pt,a4paper]{article}
\usepackage[utf8]{inputenc}
\usepackage{graphicx}
\usepackage{booktabs}
\usepackage{longtable}
\usepackage{hyperref}
\usepackage{geometry}
\usepackage{amsmath}
\usepackage{caption}
\geometry{margin=2cm}

\title{Tema 7 – Procura Sistemática de Literatura}
\author{Departamento de Física, Universidade Eduardo Mondlane}
\date{\today}

\begin{document}
\maketitle

\tableofcontents
\newpage

%-----------------------------------
\section{Introdução}
As revisões sistemáticas são uma ferramenta essencial para consolidar evidências científicas, avaliando estudos de forma objetiva e estruturada.  
A pesquisa sistemática de literatura permite identificar, selecionar e sintetizar informações relevantes, garantindo rigor metodológico e transparência.

\subsection{Objetivos da Sessão}
\begin{itemize}
    \item Compreender o conceito e importância da pesquisa sistemática de literatura.
    \item Aprender a desenvolver estratégias de pesquisa robustas.
    \item Aplicar frameworks como PICO/ST para estruturar consultas.
    \item Explorar técnicas avançadas de pesquisa em bases de dados.
\end{itemize}

%-----------------------------------
\section{Principais Características de uma Revisão Sistemática}
\begin{itemize}
    \item Estruturada, objetiva e reproduzível.
    \item Define critérios de elegibilidade claros.
    \item Utiliza estratégias de pesquisa abrangentes.
    \item Documenta e sintetiza resultados de estudos relevantes.
\end{itemize}

\begin{figure}[h!]
\centering
\includegraphics[width=0.7\textwidth]{caracteristicas_rsm.png} % Vão
\caption{Principais características de uma revisão sistemática}
\end{figure}

%-----------------------------------
\section{Introdução à Pesquisa Sistemática de Literatura}
A pesquisa sistemática é essencial para:
\begin{itemize}
    \item Garantir abrangência e exatidão na identificação de estudos.
    \item Reduzir viés de seleção de artigos.
    \item Facilitar sínteses quantitativas e qualitativas.
\end{itemize}

\begin{figure}[h!]
\centering
\includegraphics[width=0.7\textwidth]{importancia_pesquisa.png} % Vão
\caption{Importância da pesquisa sistemática de literatura}
\end{figure}

%-----------------------------------
\section{Desenvolvimento da Estratégia de Pesquisa}
\subsection{Etapas do Processo de Pesquisa}
\begin{enumerate}
    \item Definição da questão de pesquisa (ex.: PICO/ST)
    \item Identificação de palavras-chave e sinônimos
    \item Estruturação da consulta usando operadores booleanos
    \item Pesquisa em bases de dados relevantes
    \item Aplicação de filtros e critérios de elegibilidade
    \item Extração e organização dos resultados
\end{enumerate}

\begin{figure}[h!]
\centering
\includegraphics[width=0.7\textwidth]{etapas_pesquisa.png} % Vão
\caption{Etapas do processo de pesquisa sistemática}
\end{figure}

%-----------------------------------
\section{Estruturação da Estratégia de Pesquisa}
\subsection{Framework PICO/ST}
\begin{itemize}
    \item \textbf{P} – População
    \item \textbf{I} – Intervenção/Exposição
    \item \textbf{C} – Comparador
    \item \textbf{O} – Desfecho
    \item \textbf{S} – Tipo de estudo
    \item \textbf{T} – Tempo
\end{itemize}

\begin{table}[h!]
\centering
\begin{tabular}{p{2.5cm} p{10cm}}
\toprule
Elemento & Exemplo de Aplicação \\
\midrule
P & Pacientes com diabetes tipo 2 \\
I & Intervenção nutricional baseada em baixo índice glicêmico \\
C & Dieta padrão \\
O & Controle glicêmico (HbA1c) \\
S & Ensaios clínicos randomizados \\
T & Seguimento de 6 meses \\
\bottomrule
\end{tabular}
\caption{Exemplo de questão estruturada PICO/ST}
\end{table}

%-----------------------------------
\section{Técnicas de Pesquisa}
\subsection{Palavras-chave e Sinônimos}
\begin{itemize}
    \item Identificar termos principais e alternativos.
    \item Usar dicionários de sinônimos e Tesauros de bases de dados.
\end{itemize}

\subsection{Cabeçalhos de Assunto e Operadores Booleanos}
\begin{itemize}
    \item \textbf{OR:} combina sinônimos e termos relacionados
    \item \textbf{AND:} combina conceitos diferentes
    \item \textbf{NOT:} exclui termos irrelevantes
\end{itemize}

\begin{figure}[h!]
\centering
\includegraphics[width=0.6\textwidth]{operadores_booleanos.png} % Vão
\caption{Exemplo de operadores booleanos}
\end{figure}

%-----------------------------------
\subsection{Truncamento e Proximidade}
\begin{itemize}
    \item Truncamento: $*$(asterisco) para incluir variações de palavras.
    \item Proximidade: busca por termos próximos no texto.
\end{itemize}

\begin{figure}[h!]
\centering
\includegraphics[width=0.6\textwidth]{truncamento_proximidade.png} % Vão
\caption{Exemplo de truncamento e busca por proximidade}
\end{figure}

%-----------------------------------
\section{Execução da Pesquisa em Bases de Dados}
\begin{itemize}
    \item Scopus, Web of Science, PubMed, Google Scholar.
    \item Uso de filtros: idioma, ano, tipo de estudo.
    \item Coleta de artigos completos para análise.
\end{itemize}

\begin{figure}[h!]
\centering
\includegraphics[width=0.7\textwidth]{bases_dados.png} % Vão
\caption{Bases de dados e estratégias de pesquisa}
\end{figure}

%-----------------------------------
\section{Exemplo Prático de Consulta}
\begin{table}[h!]
\centering
\begin{tabular}{p{4cm} p{10cm}}
\toprule
Base de dados & Exemplo de consulta \\
\midrule
Scopus & TITLE-ABS-KEY(diabetes AND "low glycemic index" AND nutrition) \\
Web of Science & TS=(diabetes AND "low glycemic index" AND nutrition) \\
PubMed & ("diabetes"[Mesh] AND "low glycemic index"[tiab] AND nutrition[tiab]) \\
\bottomrule
\end{tabular}
\caption{Exemplo de consultas em bases de dados}
\end{table}

%-----------------------------------
\section{Exercício Prático}
\begin{itemize}
    \item Desenvolver uma estratégia de pesquisa usando PICO/ST.
    \item Aplicar operadores booleanos e truncamento.
    \item Executar a pesquisa em pelo menos uma base de dados.
    \item Coletar e organizar resultados.
\end{itemize}

\begin{figure}[h!]
\centering
\includegraphics[width=0.7\textwidth]{exercicio_pesquisa.png} % Vão
\caption{Exercício prático de pesquisa sistemática}
\end{figure}

%-----------------------------------
\section*{Referências}
\begin{itemize}
    \item Higgins JPT, Thomas J, Chandler J, et al. \textit{Cochrane Handbook for Systematic Reviews of Interventions}, 2nd Edition, 2019.
    \item Borenstein M., Hedges L., Higgins J., Rothstein H. \textit{Introduction to Meta-Analysis}, 2009.
    \item Moher D., Liberati A., Tetzlaff J., Altman DG. \textit{Preferred Reporting Items for Systematic Reviews and Meta-Analyses: The PRISMA statement}, 2009.
\end{itemize}

\end{document}
